%%
%% spring-lecture10.tex
%% 
%% Made by Alex Nelson <pqnelson@gmail.com>
%% Login   <alex@lisp>
%% 
%% Started on  2026-04-28T10:27:49-0700
%% Last update 2026-04-28T10:27:49-0700
%% 

\lecture[Jacobson ideal]

% Pierce, section 4.3

\begin{lemma}\label{lemma:4.8}
Let us write ${}_{A}A$ for treating $A$ as a left $A$-module, and
$A_{A}$ for treating $A$ as a right $A$-module. Then $\Rad({}_{A}A)=\Rad(A_{A})$.
\end{lemma}

\begin{proof}
By the left version of Proposition~\ref{prop:4.2} and Lemma~\ref{lemma:4.5},
we can take $P=\Rad({}_{A}A)\subset A_{A}$ satisfies Nakayama's lemma
for modules $\forall N\subset M\ldotp P+N=M\implies N=M$. Then using
Lemma~\ref{lemma:4.6}, $\Rad({}_{A}A)\subset\Rad(A_{A})$. In exactly
the same manner, $\Rad(A_{A})\subset\Rad({}_{A}A)$.
\end{proof}

\begin{definition}
For any algebra $A$, we have the \define{Jacobson Ideal} of $A$ be $\Jacobson(A):=\Rad(A_{A})$.
\end{definition}

\begin{proposition}\label{prop:4.9}
The Jacobson ideal of $A$ is a two-sided ideal $\Jacobson(A)\in I(A)$
such that
\begin{equation*}
\Jacobson(A)=\{x\in A\mid\forall y\in A\ldotp 1+xy\in A^{\times}\}
\end{equation*}
\end{proposition}

\begin{proof}
This follows from Lemma~\ref{lemma:4.6} and Lemma~\ref{lemma:4.8}.
\end{proof}

\begin{fact}
\begin{equation*}
\begin{split}
\Jacobson(A) &= \bigcap\{M\properideal A\mid M\mbox{ is a maximal right ideal of }A\}\\
&= \bigcap\{M\properideal A\mid M\mbox{ is a maximal left ideal of }A\}
\end{split}
\end{equation*}
\end{fact}

\begin{corollary}\label{cor:4.10}
Let $M$ be a right (or left) $A$-module.
\begin{enumerate}
\item If $1+x\in A^{\times}$ is a unit for all $x\in M$, then $M\subset\Jacobson(A)$.
If also $\Rad(A/M)=0$, then $M=\Jacobson(A)$.
\item If every $x\in M$ is nilpotent, then $M\subset\Jacobson(A)$.
\end{enumerate}
\end{corollary}

\begin{proof}
\begin{enumerate}
\item Follows from Proposition~\ref{prop:4.9} and Lemma~\ref{lemma:2.28}.
\item If $x^{n}=0$, then
\begin{equation}
(1+x)^{n}\left(\sum^{n-1}_{j=0}(-x)^{j}\right)=\left(\sum_{j}(-x)^{j}\right)(1+x)^{n}=1,
\end{equation}
so $1+x\in A^{\times}$ is a unit. The result follows from the first
claim in this corollary.\qedhere
\end{enumerate}
\end{proof}

\begin{lemma}\label{lemma:4.11}
\begin{enumerate}
\item If $\theta A\to B$ is a surjective algebra morphism, then $\theta(\Jacobson(A))\subset\Jacobson(B)$
\item $\Jacobson(A\times B)=\Jacobson(A)\times\Jacobson(B)$
\end{enumerate}
\end{lemma}

\begin{proof}
\begin{enumerate}
\item Follows from Lemma~\ref{lemma:4.1} (the image of the radical is
  contained in the raical) and Lemma~\ref{lemma:2.1} that
  $\Rad(B_{A})=\Rad(B_{B})$ which implies the result.
\item By (1), the inclusion morphisms gives us $\Jacobson(A\times B)\subset\Jacobson(A)\times\Jacobson(B)$.


For $x\in\Jacobson(A)$ and $y\in\Jacobson(B)$, by
Proposition~\ref{prop:4.9} $1_{A}+x\in A^{\times}$ and $1_{B}+y\in B^{\times}$,
so $(1_{A}+x,1_{B}+y)\in(A\times b)^{\times}$ and
\begin{subequations}
  \begin{align}
(1_{A}+x,1_{B}+y) &= (1_{A},1_{B})+(x,y)\\
&=1_{A\times B}+(x,y).
  \end{align}
\end{subequations}
Since $\Jacobson(A)\times\Jacobson(B)\in I(A\times B)$, by
Corollary~\ref{cor:4.10}~(1) we have
$\Jacobson(A)\times\Jacobson(B)\subset\Jacobson(A\times B)$ which
proves the claim.\qedhere
\end{enumerate}
\end{proof}

\begin{example}[Homework]\label{ex:4.12}
If $M$ is a semisimple $A$-module, then the Jacobson ideal for the
endomorphism algebra of $M$ vanishes: $\Jacobson(\End_{A}(M))=0$.

Further, if $M$ is finitely-generated, then $\End_{A}(M)$ is semisimple.
\end{example}

\begin{proposition}\label{prop:4.13}
If $A$ is right or left Artinian, thene there exists some $k\in\NN$
such that $\Jacobson(A)^{k}=0$ (i.e., every $x\in\Jacobson(A)$
satisfies $x^{k}=0$).
\end{proposition}

\begin{proof}
We have the chain of ideals $\Jacobson(A)\supset\Jacobson(A)^{2}\supset\Jacobson(A)^{3}\supset\dots$.
Since $A$ is Artinian, there exists a $k\in\NN$ such that
$\Jacobson(A)^{k}=\Jacobson(A)^{k+1}=\Jacobson(A)^{k+2}=\dots$.

Assume for contradiction $\Jacobson(A)^{k}\neq0$. Let
\begin{equation}
S=\{\mbox{nonzero right ideals }M\ideal A\mid M\Jacobson(A)=M\}\subset S(A)
\end{equation}
we can take $M$ the minimal element of $S$. Then since
\begin{equation}
0\neq M=M\Jacobson(A)=M\Jacobson(A)^{2}=\dots=M\Jacobson(A)^{k},
\end{equation}
there must be some $x\in M$ such that
\begin{equation}
x\Jacobson(A)^{k}\neq0.
\end{equation}
But $x\Jacobson(A)^{k}$ is a right ideal and so
\begin{equation}
x\Jacobson(A)^{k}\Jacobson(A)=x\Jacobson(A)^{k+1}=x\Jacobson(A)^{k},
\end{equation}
so $x\Jacobson(A)\in S$.

By minimality of $M$,
\begin{equation}
M=x\Jacobson(A)^{k}\subset xA\subset M,
\end{equation}
so $M$ is finitely-generated. By Corollary~\ref{cor:4.7}, $M=0$ which
gives us our contradiction.
\end{proof}

\begin{corollary}\label{cor:4.14}
For any right (resp., left) Artinian algebra $A$ and right (resp.,
left) ideal $M$ of $A$, the following are equivalent:
\begin{enumerate}
\item $M\subset\Jacobson(A)$
\item $\exists k\in\NN\ldotp M^{k}=0$
\item every element $x\in M$ is niplotent
\end{enumerate}
\end{corollary}

\begin{proof}
$(1)\implies(2)$ by Proposition~\ref{prop:4.13}.

$(2)\implies(3)$ obviously.

$(3)\implies(1)$ by Corollary~\ref{cor:4.10}~(2).
\end{proof}

\begin{proposition}\label{prop:4.15}
For any left (resp., right) Artinian algebra $A$, every Artinian
$A$-module is Noetherian.
\end{proposition}

\begin{proof}
Let us write $J=\Jacobson(A)$. Let $M$ be an arbitrary Artinian
$A$-module.

Since $A$ is Artinian, there exists
some $k\in\NN$ such that $J^{k}=0$ by Proposition~\ref{prop:4.13}. In
particular, we may take the smallest $n\leq k$ such that $MJ^{n}=0$.

We will prove the claim ($MJ^{n}=0$ implies $M$ is Noetherian) by induction on $n$.

\textsc{Case 1} ($n=0$): Then $M=0$, hence $M$ is Noetherian.

\textsc{Case 2} ($n=1$): We see $MJ=0$, which implies $M$ is an
$(A/J)$-module. Since $A/J$ is semisimple (by Corollary~\ref{cor:4.3}~(1)),
every $(A/J)$-module is semisimply by Proposition~\ref{prop:3.2}~(2).
Therefore $M_{A/J}$ (which is $M$ viewed as an $(A/J)$-module) is
Noetherian by Proposition~\ref{prop:2.26}. Since $S(M_{A})=S(M_{A/J})$,
and $S(M_{A/J})$ is Noetherian, it follows that $S(M_{A})$ is
Noetherian. Hence $M$ is Noetherian.

\textsc{Case 3} ($n>1$): Assume $MJ^{n}=0$. Let $0\neq N=MJ^{n-1}\subset M$
be a submodule of $M$. (Note $N\neq0$ by minimality of $n$.) Then $N$
is Artinian and $NJ=0$. By the $n=1$ case, $N$ is Noetherian. Also
$M/N$ is Artinian, so $(M/N)J^{n-1}=0$, so by the inductive hypothesis
$M/N$ is Noetherian. Therefore $M$ is Noetherian by Lemma~\ref{lemma:2.22}.
\end{proof}

\begin{corollary}\label{cor:4.16}
Let $A$ be a left (resp., right) Artinian algebra, let $M$ be an
$A$-module. The following are equivalent:
\begin{enumerate}
\item $M$ is Artinian
\item $M$ is Noetherian
\item $M$ is finitely-generated
\end{enumerate}
\end{corollary}

\begin{proof}
$(1)\implies(2)$ by Proposition~\ref{prop:4.15}.

$(2)\implies(3)$ by Proposition~\ref{prop:2.26}.

$(3)\implies(1)$ Write $M=\sum^{n}_{i=1}u_{i}A$. There exists a
surjective $A$-module morphism
\begin{equation}
\begin{split}
\bigoplus nA\to M\\
(x_{i})\mapsto\sum u_{i}x_{i}.
\end{split}
\end{equation}
Since $A$ is Artinian, $\bigoplus nA$ and $M$ are Artinian by
Lemma~\ref{lemma:2.22}. Hence the claim.
\end{proof}

\begin{note}
These previous few results implies that Artinian algebras are Noetherian.
\end{note}

\begin{recall}
For $\alpha\in M_{n}(\FF)$ with components $\alpha=(\alpha_{i,j})$,
that the trace is $\tr(\alpha)=\sum^{n}_{i=1}\alpha_{i,i}$ the sum of
the diagonal components. This mapping
\begin{equation}
\tr\colon M_{n}(\FF)\to\FF
\end{equation}
is $\FF$-linear.

If $\alpha$ is a nilpotent matrix, then $\tr(\alpha)=0$.
If $\lambda$ is an eigenvalue of $\alpha$, so $\alpha v=\lambda v$,
then $\alpha^{n}v=\lambda^{n}v=0$, which implies $\lambda=0$.
\end{recall}

\begin{lemma}\label{lemma:4.17}
There does not exist a set of niplotent matrices that spans
$M_{n}(\FF)$ as an $\FF$-space.
\end{lemma}

\begin{proof}
If there exists nilpotent $\alpha_{1},\dots,\alpha_{r}\in M_{n}(\FF)$
and there exists $b_{1},\dots,b_{r}\in\FF$ such that
\begin{equation}
\canMat_{1,1}=\sum^{r}_{i=1}\alpha_{i}b_{i},
\end{equation}
then
\begin{subequations}
  \begin{align}
1 &=\tr(\canMat_{1,1})\\
&=\tr\left(\sum^{r}_{i=1}\alpha_{i}b_{i}\right)\\
&=\sum^{r}_{i=1}\tr(\alpha_{i})b_{i}\\
&=0,
  \end{align}
\end{subequations}
which is a contradiction. Hence the claim.
\end{proof}

\begin{proposition}\label{prop:4.18}
Let $A$ be a finite-dimensional $\FF$-algebra, let $B$ be an
$\FF$-subspace of $A$ that is closed under multiplication and spanned
by a set of nilpotent elements (of $A$). Then $B^{k}=0$ for some $k\in\NN$.
\end{proposition}

\begin{proof}
There are two simplifications which occur for this proof.
  
First simplification: we may assume that $\FF$ is algebraically closed. In fact, let
$x_{1}$, \dots, $x_{r}$ be an $\FF$-basis for $A$ (as a vector space)
such that
\begin{equation}
x_{j}x_{k}=\sum^{r}_{i=1}c^{i}_{jk}x_{i}
\end{equation}
for some coefficients $c^{i}_{jk}\in\FF$. Then for an
$\closure{\FF}$-algebra $A'_{\FF}=\bigoplus_{i} x_{i}\closure{\FF}$ and
multiplication is given by the structure constants $c^{i}_{jk}$. Then
$A$ is a subalgebra of $A'_{\FF}$. The $B'=B\closure{\FF}$ satisfies
the same condition as $B$ and $(B')^{k}=0$ implies $B^{k}=0$.

Second simplification: $B$ is an ideal of $A$. This is because $B$ is
an ideal of $B+1_{A}\FF$, and $B+1_{A}\FF$ is an ideal of $A$.

\textsc{Claim:} if $A$ is semisimple, then $B=0$.

We will stipulate the claim holds, prove it is sufficient to finish
our proof, then we will prove the claim. If $A$ is not
semisimple, then we can replace $A$ by $A/\Jacobson(A)$ (which
\emph{is} semisimple since $\dim_{\FF}(A)$ is finite which implies $A$
is right Artinian and then we use Corollary~\ref{cor:4.3}~(1)) and we replace $B$ by the ideal
$(B+\Jacobson(A))/\Jacobson(A)$ of $A/\Jacobson(A)$. If the claim
holds, then $B+\Jacobson(A)=\Jacobson(A)$, which implies
$B\subset\Jacobson(A)$ and so $B^{k}=0$ for some $k$ by
Proposition~\ref{prop:4.13}.

Let us now prove the claim. By Corollary~\ref{cor:3.22},
\begin{equation}
A=\prod^{t}_{i=1}A_{i}
\end{equation}
where each $A_{i}\iso M_{n_{i}}(\FF)$. For the projection morphism
\begin{equation}
\pi_{i}\colon A\to A_{i},
\end{equation}
we see $\pi_{i}(B)$ is an ideal of the simple algebra $A_{i}$. This
means $\pi_{i}(B)=0$ or $\pi_{i}(B)=A_{i}$. But $\pi_{i}(B)=A_{i}$ is
impossible because $\pi_{i}(B)$ is nilpotent. Therefore $\pi_{i}(B)=0$
for every $i$. Then
\begin{equation}
B\subset\bigcap^{t}_{i=1}\ker(\pi_{i})=0,
\end{equation}
which proves the claim.
\end{proof}